% Generated by D:/tools/scratch_qdd/mmrec/render.py from data_f.py - edit the data, not this file.
\subsection{F. 과정과 규칙 --- 기록, 자체 검사, 영상, 절약, 논문 갱신, 이름}\label{rec:F}
\noindent\textbf{한 줄:} 모든 결정과 실수를 MD·기록책에 남기고(함정 P01--P123), 모든 실험은 사전 등록 + 시작 전·도중·끝 자체 검사, 폐루프 판은 영상 필수, 유료는 명시 승인 뒤. 논문은 가설 기반 완성판 + 빨간 줄 마인드맵 두 판을 4시간마다 갱신. 논문 이름은 PaceNotes(저장소·코드 이름 harvest 유지).

\paragraph{F1. 모든 것을 MD로 --- 같은 실수를 반복하지 않게}
{\small 09-23 \textasciitilde{} 09-24 04:58 KST}
\begin{itemize}
  \item \textbf{사용자} (ul 12·16, 09-23 첫 세션): \intent{초안 작성 버전도 MD로 다 기록한다(이때는 이렇게 했었다). 같은 실수를 반복하지 않기 위해서다.}
  \item \textbf{사용자} (ul 12·16, 09-23 첫 세션): \intent{사용자가 의도하고 말한 것은 빨간색, 나머지는 검정색.}
  \item \textbf{사용자} (ul 20, 09-24 04:50 KST 경): \intent{꼭 MD를 쓰면서 진행한다. MD에 기록해서 같은 실수를 반복하지 않는다.}
  \item \textbf{그래서(결정):} CLAUDE.md(작업 규칙)·\texttt{docs/\allowbreak{}user-\allowbreak{}log.\allowbreak{}md}(사용자 발언 전부)·\texttt{docs/\allowbreak{}draft-\allowbreak{}log.\allowbreak{}md}(판마다 한 일·틀린 것·교훈)·정본 \texttt{docs/\allowbreak{}design/\allowbreak{}00-\allowbreak{}interfaces.\allowbreak{}md}(뒤 절이 앞 절보다 우선). 마인드맵의 빨간 줄은 user-log 글자와 일치해야 하고 \texttt{tools/\allowbreak{}intent\_\allowbreak{}check.\allowbreak{}py}로 확인한다.
\end{itemize}

\paragraph{F2. 논문은 두 판 --- 논문형 본문과 빨간 줄 마인드맵}
{\small 09-24 04:58 KST \textasciitilde{} 09-27}
\begin{itemize}
  \item \textbf{사용자} (ul 21, 09-24 04:58 KST): \intent{Overleaf는 절대 논문이 아니다. 마인드맵 그리듯 초안을 잡는 느낌이다.}
  \item \textbf{사용자} (ul 26, 09-24 12:32 KST): \intent{내용이 얼마나 많고 적음은 따지지 말고 지금 한 것까지만 해서 글을 좀 다듬어 줘봐 지금 있는 내용에서만 글을 다듬어 달라는 얘기야 막 어, 추가로 뭘 넣거나 이런다는 얘기가 아니라}
  \item \textbf{사용자} (ul 26, 09-24 12:32 KST): \intent{실제 논문처럼 글을 써주면 돼.}
  \item \textbf{그래서(결정):} 본문 \texttt{paper/\allowbreak{}main.\allowbreak{}tex} + \texttt{sec/\allowbreak{}*.\allowbreak{}tex} = 논문형 한국어 글(있는 내용만), 마인드맵 \texttt{paper/\allowbreak{}mindmap.\allowbreak{}tex} + \texttt{mindmap/\allowbreak{}sec/\allowbreak{}*.\allowbreak{}tex} = 빨간 줄(사용자 의도)을 보관하며 계속 갱신.
  \item \textbf{사용자} (ul 63, 09-25 03:18 KST): \intent{지금까지 된 거 어떤식으로 할건지 해서 실제 논문형식으로 오버리프 다시 확 다듬어줘 아직 확실한 자료가 아닌건 임시로 넣어두고 하는 것도 괜찮아}
  \item \textbf{사용자} (ul 88, 09-26 03:02 KST): \intent{지금의 가설을 기반으로 지금 아직 진행안된거 기바ㅏㄴ으로도 해서 제일 이상적인걸로 정리최대해서 논문 ㅈㅇ리해서 지금바로 줘봐보바 당연히 나중에 계속 업뎃 동일하게해주고}
  \item \textbf{그래서(결정):} 정본 §85·CLAUDE.md(13f3495, 03:02): 논문은 가설 기반 이상적 완성판 --- 안 한 실험은 가설·등록 계획·예상으로 `예정' 표시(주황), 숫자 날조 금지. 첫 판 745af38(03:19).
  \item \textbf{참고:} ul 113(B5)으로 확정 전 가설도 그 자리에서 논문(주황)·마인드맵에 올린다. 
  \item \textbf{사용자} (ul 132, 09-27 02:0x KST (통제자 경유)): \intent{자 논문을 우리 2개 놓고 있잖아? 그 한개는 내가말한말 있는거에는 엄청다 적는걸로 하자 이렇게 해봤는데 ㅇ래서 이렇게 하기로 했다 이런식의 상세한것까지. 단 원본 논문은 이제 진짜 다듬는걸 한번 실행을 해보자ㅣ}
  \item \textbf{사용자} (ul 133, 09-27 02:0x KST (통제자 경유)): \intent{그 피규어들도 전부 다 돌면서 새로 만드는거나 삭제, 수정 다 해줘봐봐 메인 피규어도 마찬가지}
  \item \textbf{그래서(결정):} 두 판을 나눔: 본문은 다듬기 판(14930f7 --- 이야기 한 줄기, 실험 장부성 내용은 부록·마인드맵으로, 50 $\rightarrow$ 14쪽)과 그림 정리(논문 에이전트), 이 마인드맵은 상세 결정 이력(5절, 17395b3)이고 제목을 `PaceNotes --- 연구 기록·마인드맵 (상세 결정 이력)'으로.
\end{itemize}

\paragraph{F3. 논문 갱신 주기: 1시간 $\rightarrow$ 4시간, 큰 결과는 수시}
{\small 09-25 18:25 KST \textasciitilde{}}
\begin{itemize}
  \item \textbf{사용자} (ul 68·69, 09-25 18:25--18:28 KST): \intent{지속적으로 오버리프 업데이트 하면서 하는중?}
  \item \textbf{사용자} (ul 68·69, 09-25 18:25--18:28 KST): \intent{1시간마다 논문 업데이트하기로해 앞으로는 꼭}
  \item \textbf{사용자} (ul 73, 09-25 23:04 KST): \intent{논문 업뎃주기는 4시간으로 하자}
  \item \textbf{그래서(결정):} CLAUDE.md(93c68d7, 585786d): 4시간 주기(매 4시간 :17), 큰 판정은 주기와 별도 반영. 정기 갱신 예: cd5210f(09-26 01:07), 54b9205(08:57), 5ea8585(12:56), 858ac6c(18:50), 9d12a4a(20:52), 1d3c86c(09-27 00:51); 수시 갱신 054f250(H-L), 35c9fbb(L8).
\end{itemize}

\paragraph{F4. 조사 규칙과 보고 시각}
{\small 09-23 \textasciitilde{} 09-25 23:03 KST}
\begin{itemize}
  \item \textbf{사용자} (ul 14, 09-23): \intent{신뢰도 낮은 논문과 깃 저장소는 최대한 쓰지 않는다. 쓰느니만 못하다.}
  \item \textbf{사용자} (ul 70, 09-25 19:03 KST): \intent{문제가 있다면 은 만약에 성능이 너무 낮다면 은 기존에 최근 1년 반 이후 논문에서는 그러한 성능들을 비슷한 논문을 찾아가지고 그러한 성능들을 어떻게 올렸는지 찾아봐서 적응을 해도}
  \item \textbf{그래서(결정):} 문헌은 1년 반 이내(큰 주제의 기초 문헌은 `기간 밖, 기초' 표시), 학회·인용·별로 신뢰도 확인. 성능이 낮으면 선행 연구의 개선법을 새 판·새 사전 등록으로 적용(문턱을 결과 뒤에 바꾸지 않음) --- 예: 준수율 0.8 조사(D5), E-SR1e의 L1 대조 안내(D9).
  \item \textbf{사용자} (ul 72, 09-25 23:03 KST): \intent{한국시간으로 해줘 앞으론}
  \item \textbf{그래서(결정):} 사용자 보고는 KST, 저장소 기록은 UTC(d734d00). 이 5절은 읽는 사람을 위해 KST로 적었다.
\end{itemize}

\paragraph{F5. 실험마다 자체 검사}
{\small 09-26 02:41 KST \textasciitilde{}}
\begin{itemize}
  \item \textbf{사용자} (ul 87, 02:45 KST): \intent{매번 자체검사를 해야해 알겠지?}
  \item \textbf{그래서(결정):} CLAUDE.md(ed478e4): 모든 실험(유료·GPU)은 (1) 시작 전 --- 바꾸는 설계 결정, 이 표본으로 가를 수 있나, 더 싼 사전 실행을 사전 등록에; (2) 도중 --- 관문 결함(버그·무효 > 10 \%·교란·비용이나 시간 2배·관문 실패)이면 멈추고 재설계 $\rightarrow$ 등록 변경을 커밋한 뒤 재개; (3) 끝 --- 검증 명령으로 확인 뒤 보고.
  \item \textbf{결과:} 실제 적용 예: 탐침 재범위(A3), E-ACC 변경 1--6(A6·A7), Astra 단독 변경 1--3(A8·A9), E-SR1c 변경 1·2(D7), E-PT 변경 1(B7), E-STRIP8 변경 1(B9). 규칙 이탈도 결과 문서에 적는다 --- heredoc·\texttt{python -\allowbreak{}}·\texttt{python -\allowbreak{}c}를 쓴 일이 여러 에이전트에서 되풀이됨(함정 P03).
  \item \textbf{참고:} 유료 쪽 `승인 불필요' 부분은 ul 114가 대체했다(A10). 자체 검사 의무는 그대로.
\end{itemize}

\paragraph{F6. 결과·실수를 책처럼 --- 연구 기록책}
{\small 09-26 03:34 -- 03:59 KST}
\begin{itemize}
  \item \textbf{사용자} (ul 89, 03:34 KST): \intent{결과들은 모두 효율적으로 정보 저장해둬 나중에 같은씰수안하게 동일 md를 써도되고 효율적으로 분산해서 책처럼 기록 남겨둬두되고}
  \item \textbf{사용자} (ul 89, 03:34 KST): \intent{업데이트 기록 필수임}
  \item \textbf{그래서(결정):} \texttt{docs/\allowbreak{}book/\allowbreak{}}(e15270a, 03:40): 01 실험 장부·02 함정과 예방 규칙·03 결정 연표·04 데이터/코드/파드 지도·05 비용 장부; 06 갱신 기록 필수(4117851, 03:59). 결과·실수가 생기면 같은 커밋에 책도 갱신, 새 작업 전에 함정 장을 먼저 읽는다.
  \item \textbf{결과:} 지금 함정 장은 P123까지 있다. 예: P80(경로 없이 커밋해 남의 파일을 쓸어 담음 $\rightarrow$ 늘 경로 지정 커밋), P101(불리한 기준선), P108(크레딧 소진을 빈 답으로 셈), P113(외삽 예상 상태), P122(IK 잠김 막힘 행), P123(vLLM 환경·병합 샤드 색인).
\end{itemize}

\paragraph{F7. 영상 규칙}
{\small 09-26 19:21 KST \textasciitilde{}}
\begin{itemize}
  \item \textbf{사용자} (ul 131, 09-26 19:2x KST (통제자 경유)): \intent{헐 그거. 영상들도. 무조건 다 데이타에 저장해둬}
  \item \textbf{사용자} (ul 131, 09-26 19:2x KST (통제자 경유)): \intent{필수야 심서 아스트라가 조금이라도 돌릴때 영상 저장.}
  \item \textbf{사용자} (ul 131, 09-26 19:2x KST (통제자 경유)): \intent{노노 그냥 아까 아스트라 단독만 일단}
  \item \textbf{그래서(결정):} Astra 단독 실제 편 7개 영상(편마다 머리 | 오른손목 연속 영상 + 모델 시점 영상)을 파드 \texttt{/\allowbreak{}data/\allowbreak{}harvest/\allowbreak{}videos/\allowbreak{}}에 저장(4d54506). 이후 폐루프 시험의 등록·결과는 영상을 필수로 적는다: E-OpenVLM 대리(\texttt{videos/\allowbreak{}open\_\allowbreak{}vlm\_\allowbreak{}proxy/\allowbreak{}}), E-TEACH-L8(\texttt{videos/\allowbreak{}teach\_\allowbreak{}l8/\allowbreak{}}), E-PT·E-STRIP8 등록(폐루프 영상).
\end{itemize}

\paragraph{F8. 절약과 작업 위치}
{\small 09-24 17:30 KST \textasciitilde{}}
\begin{itemize}
  \item \textbf{사용자} (ul 43, 09-24 17:30 KST): \intent{넣는 건데 2번, 3번, 3번은 일단 비용 한도 없이 그냥 할 거니까 진행해 주되 최대한 비용 문제 없게 진행을 해주고 2번 경우에는 어, 사용해도 돼 1번의 경우에는 일단 후순위로 미루자 내가 나중에 줄게}
  \item \textbf{사용자} (ul 45·51, 09-24 18:24 · 21:20 KST): \intent{H200 서버 사용할 때 무조건 데이터 경로에다가 사용해줘야 돼.}
  \item \textbf{사용자} (ul 45·51, 09-24 18:24 · 21:20 KST): \intent{필요한 큰거나 기본은 전부다 data 경로인건 알지? 쿠베의 h200의}
  \item \textbf{사용자} (ul 45·51, 09-24 18:24 · 21:20 KST): \intent{폴더 이름은 harvest 여야지}
  \item \textbf{그래서(결정):} 유료 한도 약 10만 원(ul 76) $\rightarrow$ 실험별 상한·80 \% 정지 $\rightarrow$ 이제 명시 승인 뒤에만(ul 114, A10). 무료 모델로 먼저 배선·버그를 잡고 유료는 가를 수 있는 표본만(탐침 재범위, Astra 단독 단계식). 파드 파일은 모두 \texttt{/\allowbreak{}data/\allowbreak{}harvest/\allowbreak{}} 아래, 로컬 임시 파일은 \texttt{D:\textbackslash{}tools\textbackslash{}scratch\_\allowbreak{}qdd}(C 드라이브 금지).
\end{itemize}

\paragraph{F9. 논문 이름 PaceNotes}
{\small 09-27 01:46 KST (커밋)}
\begin{itemize}
  \item \textbf{사용자} (ul 129, 09-27 01:5x KST): \intent{우리 이거 일단 논문을 정리를 할건데 이 논문의 이름을 뭐라할까? 내가 이논문이름을 원래 harvest로 했던 이유는 seed라고 그래스프잰뿌려리고 그리고 mill 그중에 장애물이나 잡기쉬운걸로 하고 하베스트 그중에 뽈라서 실행에 올린다. 그다음에 그걸 학습해서 bake한다는 이게 그냥 옛날아이디어야 이아이디어는 무시해두되는데 이런식으로 harvest를 가져가도 좋고 다른 아이디어 있으면 그걸로 한번 해줘봐도돼 후보를 만들어보자 재치있음 좋겠어. 개인적으로 나는 이게 엄청난 혁신이 될 수도 있겠단 생각을 해서 그래}
  \item \textbf{사용자} (ul 129, 09-27 01:5x KST): \intent{PaceNotes가 제일 잘 맞는듯}
  \item \textbf{그래서(결정):} user-log 129(f04b41d): 제목 `PaceNotes: Intent-Sharing Hierarchies for Fast, Verifiable Robot Manipulation' --- 랠리 코드라이버가 앞 구간 노트를 불러 주고(상위 의도 선언), 드라이버가 실시간 조향·제동 타이밍을 잡으며(VLA 세부 조종·잡기 타이밍), 코드라이버가 위치를 확인(상위 검증)하는 비유(§96 보충 2와 같은 구조). 저장소·코드 이름 harvest는 유지.
  \item \textbf{참고:} user-log 129의 시각 표기(16:5x UTC = 01:5x KST)가 커밋 시각(01:46 KST)보다 늦다 --- 기록 그대로 둔다.
\end{itemize}

\paragraph{F10. 논문도 상위 중심, 결합은 §96 보충 5로}
{\small 09-27 03:2x KST}
\begin{itemize}
  \item \textbf{사용자} (ul 140, 03:2x KST): \intent{논문도 이런식으로 적는걸로 하자 일단 실험들은 전부 상위를 중심으로 해야하는 것은 맞구 말이야.}
  \item \textbf{그래서(결정):} 논문 da3b98f: 근거 5절을 `상위 계획기(주 줄기)'와 `2단계(조건부)'로 나눔, 3.4절·Fig. 1·모델 그림을 §96 보충 5로, 커밋 안 된 결과는 `진행 중'으로만.
\end{itemize}

\paragraph{F11. 버리기 전 재검증 + 살아난 안은 공정 비교 한 번 더}
{\small 09-27 03:4x KST}
\begin{itemize}
  \item \textbf{사용자} (ul 143, 03:4x KST): \intent{결정 버리기전에 재검증해서 다시 살아났다고 기존거 ㄹ또 버리지말고 두개 비교를 공정하게 한번더 하는과정이 있어야해 알겠지?}
  \item \textbf{그래서(결정):} 협업 보드 규칙·메모리(b02ff17): 재검증으로 살아난 안과 기존 안은 같은 데이터·학습량·평가 세트·시드·판정 규칙으로 사전 등록한 짝 비교를 한 번 더 한 뒤에만 선택(E-STRIP8 손 정보, T3 재검증, 작전 T). 논문 6절 절차 규칙에 반영.
\end{itemize}
